% Kan + Sites/Sheaves (Ch 5-6)
% Lines 763-1065 of main.tex
% This file is for agent editing — will be merged back

\chapter{Kan Operations and Composition}
\label{ch:kan}

\section{Open Boxes and Fillers}

The \emph{Kan condition} is the defining property of cubical types:
every open box has a filler.

\begin{definition}[Open Box]
An \emph{open $n$-box in direction $i$ and side $b$} in a type $A$
consists of:
\begin{itemize}
  \item A \emph{bottom face}: $u : A$ (the base of the box).
  \item Faces for all directions $j \neq i$ and sides $c \in \{0,1\}$:
    $\alpha_{j,c} : \interval \to A$ with boundary conditions matching
    $u$ and each other on shared edges.
  \item The face $(i, b)$ is \emph{missing} --- this is the ``open''
    side.
\end{itemize}
\end{definition}

\begin{definition}[Kan Filler]
The \emph{Kan filler} $\hfilling^i_b(\alpha)$ is a term
$\hfilling^i_b(\alpha) : \interval \to A$ that:
\begin{enumerate}
  \item Agrees with $u$ at $i = b$:
    $\hfilling^i_b(\alpha)(b) = u$.
  \item Agrees with each $\alpha_{j,c}$ at $j = c$:
    $\hfilling^i_b(\alpha)(j = c) = \alpha_{j,c}$.
  \item The missing face $i = 1-b$ is the \emph{lid} of the box.
\end{enumerate}
\end{definition}

\section{Composition}

\emph{Composition} is the operation that composes two paths:

\begin{definition}[Homogeneous Composition]
Given $A : \Type$, $a : A$, and a partial path
$\varphi \vdash u : A$ with $u[\varphi] = a$:
\[
  \comp^i(A, [\varphi \mapsto u], a) : A
\]
This is the ``lid'' of the Kan filling of the open box with base $a$
and sides $u$.
\end{definition}

For program equivalence, composition means: if we can transform
$f$ into $h$ (by one axiom application) and $h$ into $g$ (by
another), we can compose these transformations to get $f \equiv g$.

\section{Composition for Multi-Step Proofs}

\begin{example}[Three-Step Equivalence]
Consider programs:
\begin{align}
  f &= \mathsf{sum}(\mathsf{sorted}(\mathsf{reversed}(x))) \\
  g &= \mathsf{sum}(x)
\end{align}
The equivalence proof is a three-step path:
\begin{align}
  p_1 &: \Path(f,\; \mathsf{sum}(\mathsf{sorted}(x)))
    && \text{(by $\mathsf{reversed}$ involution + $\mathsf{sorted}$
    perm-invariance)} \\
  p_2 &: \Path(\mathsf{sum}(\mathsf{sorted}(x)),\; \mathsf{sum}(x))
    && \text{(by $\mathsf{sum}$ perm-invariance)}
\end{align}
The composite $p_1 \cdot p_2 : \Path(f, g)$ is obtained by Kan
composition.  In Z3, the equational axioms achieve this automatically
--- Z3's E-matching engine chains axiom applications.
\end{example}


% ══════════════════════════════════════════════════════════
\part{Cohomology Theory}
% ══════════════════════════════════════════════════════════

\chapter{Sites, Presheaves, and Sheaves}
\label{ch:sheaves}

\section{The Type Site}

\begin{definition}[Type Site $\site$]
\label{def:type-site}
The \emph{type site} $\site$ is the category whose:
\begin{itemize}
  \item \textbf{Objects} are \emph{type fiber assignments}
    $U = (\tau_1, \ldots, \tau_n) \in \mathsf{TypeTag}^n$, specifying
    the type of each parameter.
  \item \textbf{Morphisms} $U \to V$ are \emph{type widenings}:
    $(U \to V)$ iff $\tau_i^U$ is a subtype of $\tau_i^V$ for all $i$.
  \item \textbf{Covering families}: for each object $U$, the covering
    $\{U_\tau\}_{\tau \in T}$ consists of all single-fiber restrictions.
\end{itemize}
\end{definition}

\begin{example}[Site for a Unary Function]
For $f(x)$ where $x$ can be $\mathsf{int}$, $\mathsf{str}$, or
$\mathsf{list}$, the site $\site$ has objects $U_\mathsf{int}$,
$U_\mathsf{str}$, $U_\mathsf{list}$, and a ``top'' object
$U_\mathsf{any} = U_\mathsf{int} \cup U_\mathsf{str} \cup
U_\mathsf{list}$.

Duck type inference may determine that $x$ is always an $\mathsf{int}$
(e.g., it's used in arithmetic).  In that case, the site has a single
object $U_\mathsf{int}$.
\end{example}

\section{Value Presheaves}

\begin{definition}[Value Presheaf]
\label{def:value-presheaf}
Given a program $f$ with $n$ parameters, the \emph{value presheaf}
$\Sem_f : \site^{\mathrm{op}} \to \Set$ assigns:
\begin{itemize}
  \item To each fiber assignment $U$: the Z3 term
    $\Sem_f(U) = \den{f}|_U$ (the denotation of $f$ restricted to
    parameters of types $U$).
  \item To each morphism $\sigma : U \to V$: the restriction
    $\Sem_f(\sigma) : \Sem_f(V) \to \Sem_f(U)$ (specializing from
    a wider type to a narrower one).
\end{itemize}
\end{definition}

\begin{remark}
The restriction $\Sem_f(\sigma)$ is implemented in Z3 by adding
constraints: $\den{f}|_U$ is $\den{f}$ with the additional constraint
$\mathsf{tag}(p_i) = \tau_i$ for each parameter $p_i$.  This is
the computational content of the restriction map.
\end{remark}

\section{The Sheaf Condition}

\begin{definition}[Sheaf Condition]
A presheaf $\presheaf$ on $\site$ satisfies the \emph{sheaf condition}
(or \emph{descent condition}) if: for every covering
$\{U_i \to U\}$ and family of sections $s_i \in \presheaf(U_i)$
that agree on overlaps ($\restr{s_i}{U_{ij}} = \restr{s_j}{U_{ij}}$),
there exists a \emph{unique} section $s \in \presheaf(U)$ with
$\restr{s}{U_i} = s_i$.
\end{definition}

\begin{proposition}
The value presheaf $\Sem_f$ is a sheaf on the type site.
\end{proposition}

\begin{proof}
Given sections on each fiber that agree on overlaps, the unique
global section is the Z3 term $\den{f}$ without any fiber restrictions.
The restriction to each fiber recovers the fiber-specific section
(by adding the tag constraint).

Uniqueness follows from the determinism of the Z3 model: a term
is uniquely determined by its value on all inputs, and the fiber
decomposition covers all inputs.
\end{proof}

\section{The Isomorphism Sheaf}

\begin{definition}[Isomorphism Sheaf]
\label{def:iso-sheaf}
Given two value presheaves $\Sem_f$ and $\Sem_g$, the
\emph{isomorphism sheaf} $\Iso(\Sem_f, \Sem_g)$ assigns:
\[
  \Iso(\Sem_f, \Sem_g)(U) = \begin{cases}
    1 & \text{if } \forall \vec{p} \in U.\;
      \den{f}(\vec{p}) = \den{g}(\vec{p}) \\
    0 & \text{if } \exists \vec{p} \in U.\;
      \den{f}(\vec{p}) \neq \den{g}(\vec{p})
  \end{cases}
\]
(A truth-valued sheaf indicating whether $f$ and $g$ agree on $U$.)
\end{definition}

The global sections of $\Iso(\Sem_f, \Sem_g)$ are:
\[
  \Gamma(\Iso(\Sem_f, \Sem_g)) = \begin{cases}
    1 & \text{if $f \equiv g$ globally} \\
    0 & \text{otherwise}
  \end{cases}
\]

Computing $\Gamma$ directly (asking Z3 for the full unconstrained
query) is one approach.  The \v{C}ech cohomology approach decomposes
this into local checks and gluing.


\chapter{The \v{C}ech Complex}
\label{ch:cech}

\section{Cochain Groups}

\begin{definition}[\v{C}ech Cochain Groups]
Given a covering $\covers = \{U_i\}_{i \in I}$ of the parameter space,
the \v{C}ech cochain groups of the isomorphism sheaf are:
\begin{align}
  C^0(\covers, \Iso) &= \prod_{i \in I} \Iso(U_i) \\
  C^1(\covers, \Iso) &= \prod_{(i,j) \in I^{(2)}} \Iso(U_i \cap U_j) \\
  C^2(\covers, \Iso) &= \prod_{(i,j,k) \in I^{(3)}}
    \Iso(U_i \cap U_j \cap U_k)
\end{align}
where $I^{(2)} = \{(i,j) : i < j\}$ and
$I^{(3)} = \{(i,j,k) : i < j < k\}$ are the sets of ordered pairs
and triples of covering indices.
\end{definition}

A 0-cochain $\varphi \in C^0$ is a collection of local
equivalence judgments --- one per fiber.  A 1-cochain
$\psi \in C^1$ is a collection of ``transition functions'' on
pairwise overlaps.

\section{Coboundary Operators}

\begin{definition}[Coboundary $\delta^0 : C^0 \to C^1$]
\label{def:delta0}
For $\varphi = (\varphi_i)_{i \in I} \in C^0$:
\[
  (\delta^0 \varphi)_{ij} = \restr{\varphi_j}{U_{ij}} \circ
    \restr{\varphi_i}{U_{ij}}^{-1}
\]
This measures the ``twist'' between two local equivalences on their
overlap.  If both are isomorphisms (local equivalence holds on both
fibers), the transition function is the identity.
\end{definition}

\begin{definition}[Coboundary $\delta^1 : C^1 \to C^2$]
For $\psi = (\psi_{ij})_{(i,j) \in I^{(2)}} \in C^1$:
\[
  (\delta^1 \psi)_{ijk} = \restr{\psi_{jk}}{U_{ijk}} \circ
    \restr{\psi_{ij}}{U_{ijk}}^{-1} \circ \restr{\psi_{ik}}{U_{ijk}}
\]
This is the \emph{cocycle condition}: the transition functions around
every triple overlap must compose to the identity.
\end{definition}

\begin{lemma}[$\delta^1 \circ \delta^0 = 0$]
The composition of coboundary operators is zero:
$\delta^1(\delta^0(\varphi)) = 0$ for all $\varphi \in C^0$.
\end{lemma}

\begin{proof}
Direct computation: expanding $\delta^1 \circ \delta^0$ on a triple
overlap $U_{ijk}$, all terms cancel pairwise by the group structure
of the isomorphisms.
\end{proof}

\section{Cohomology Groups}

\begin{definition}[\v{C}ech Cohomology Groups]
\begin{align}
  \Hzero(\covers, \Iso) &= \ker(\delta^0) \\
  \Hone(\covers, \Iso) &= \ker(\delta^1) / \mathrm{im}(\delta^0)
\end{align}
\end{definition}

\begin{theorem}[The Descent Theorem]
\label{thm:descent-full}
$\Sem_f \iso \Sem_g$ (global equivalence) if and only if:
\begin{enumerate}
  \item $\Hzero(\covers, \Iso) \neq 0$ (there exist local equivalences), and
  \item $\Hone(\covers, \Iso) = 0$ (local equivalences glue to a
    global one).
\end{enumerate}
\end{theorem}

\begin{proof}
$(\Rightarrow)$ If $\Sem_f \iso \Sem_g$ globally, restricting to
each fiber gives local equivalences that trivially glue (the
transition functions are all identities).

$(\Leftarrow)$ If local equivalences exist and $\Hone = 0$, every
1-cocycle is a coboundary.  The transition functions
$(\delta^0 \varphi)_{ij}$ are identities (since each $\varphi_i$ is
an isomorphism).  The gluing axiom of the sheaf gives the unique
global section extending the local equivalences.
\end{proof}

\section{Computation over GF(2)}

In our setting, the coefficient sheaf takes values in $\{0, 1\}$
(isomorphism or not).  The \v{C}ech complex reduces to a chain
complex over $\mathrm{GF}(2)$ (the field with two elements):
\[
  \mathrm{GF}(2)^{|I|} \xrightarrow{\delta^0}
  \mathrm{GF}(2)^{|I^{(2)}|} \xrightarrow{\delta^1}
  \mathrm{GF}(2)^{|I^{(3)}|}
\]

The rank of $\Hone$ over $\mathrm{GF}(2)$ is computed by:
\[
  \mathrm{rank}(\Hone) = \mathrm{rank}(\ker(\delta^1)) -
    \mathrm{rank}(\mathrm{im}(\delta^0))
\]

This is a finite linear algebra computation, decidable in
$O(n^3)$ where $n = |I^{(2)}|$.

The implementation in \texttt{cohomology.py} uses
$\mathsf{compute\_h1\_rank\_gf2}$, which performs row reduction
on the coboundary matrices.


